\section{Discussion}

\label{sec:discussion}
% ══════════════════════════════════════════════════════════════════════════════

\subsection{Lessons Learned}

\paragraph{Participation requires simplicity.} A mobile swipe-to-vote
interface increased participation 3.2x over the initial web-only version.

\paragraph{Reputation decay is essential.} Without decay, early
contributors accumulate permanent governance power. Annual 10--20\% decay
ensures active contributor influence.

\paragraph{Review staking works.} Staking improved quality and speed,
but stake sizes require careful calibration.

\paragraph{Geographic diversity requires active measures.} Despite QV's
properties, organic participation skewed toward wealthy nations. Targeted
outreach and regional ambassadors were necessary.

\subsection{Limitations}

\begin{enumerate}
  \item \textbf{Scale}: 847 active voters is small relative to national
    agencies. Governance quality at larger scale is untested.
  \item \textbf{Token volatility}: \ZOO{} price fluctuations create
    budget uncertainty. Stablecoin treasury management mitigates this.
  \item \textbf{Expertise concentration}: Specialized areas may lack
    sufficient qualified reviewers.
  \item \textbf{Regulatory status}: DAO-funded research IP rights and
    liability vary across jurisdictions.
  \item \textbf{Centralization risk}: Research Council and Treasury
    Committee hold significant power despite decentralized design.
\end{enumerate}


% ══════════════════════════════════════════════════════════════════════════════
